\section{Catholic Doctrinal Reception after the Fathers}

\subsection{Trent: the doctrine is not built on one proof-text}

The Council of Trent's thirteenth session teaches the true, real, and
substantial presence of Christ in the Eucharist, the conversion aptly called
transubstantiation, and the worship due to Christ present. Its argument is
ecclesial and sacramental: Christ's institution words, the Church's perennial
reception, and the fittingness of sacramental presence work together. John 6
belongs to the scriptural horizon, but Trent does not make the Catholic doctrine
stand or fall on whether every verse of the discourse is a direct technical
description of sacramental mode.

That distinction matters apologetically. A Catholic need not prove
transubstantiation from the lexical force of \emph{tr\=ogein}. John 6 establishes
the life-giving flesh, true food and drink, communion, and resurrection horizon.
The institution narratives identify bread and wine with Christ's Body and
Blood. Apostolic practice and 1 Corinthians supply ecclesial and judgment
contexts. Doctrinal definition receives the whole.

\subsection{Vatican II and papal teaching}

\emph{Sacrosanctum Concilium} 47 describes the Eucharistic sacrifice as the
memorial in which the sacrifice of the Cross is perpetuated and the Church is
joined to Christ's saving work.

The \emph{Catechism of the Catholic Church} 1322--1323 begins its treatment of
the Eucharist by placing participation in the Lord's sacrifice within Christian
initiation and by gathering sacrifice, memorial, communion, unity, charity,
grace, and eschatological pledge. That placement is a doctrinal reception of
the whole mystery, not evidence for the intention of a patristic author.

Paul VI's \emph{Mysterium fidei} guards the objective presence and insists that
symbolic and effective meanings do not exclude substantial conversion. John
Paul II's \emph{Ecclesia de Eucharistia} places Eucharist at the heart of the
Church and repeatedly joins sacrifice, presence, communion, and eschatological
hope. Benedict XVI's \emph{Sacramentum caritatis} presents the Eucharist as a
mystery believed, celebrated, and lived. These documents are doctrinal
reception, not evidence for what a second-century Father thought.

The controlling Catholic judgment is therefore both stronger and more careful
than the slogan ``John 6 is obviously literal.'' Christ truly gives himself;
the Church receives the whole Christ---Body, Blood, soul, and divinity---under
either species; the mode is sacramental, not cannibalistic; faith and charity
govern fruitful communion; the gift tends toward ecclesial unity and
resurrection.

